\section{Task graphs and the value of decomposition}\label{sec:C21:decomposition}
A task decomposition is a directed graph whose nodes represent work products and whose edges represent prerequisites. A node should specify a deliverable, evidence requirements, permitted tools, budget, and acceptance condition. ``Research the topic'' is too underspecified to support reliable integration. ``Identify contradictory primary results under the declared cutoff, returning source identity, claim, experimental population, and limitation'' defines an auditable information product. The distinction matters because integration quality depends on what the delegate was asked to preserve.

Let $G=(V,E)$ be an acyclic work graph, $w_v\geq0$ the execution time of node $v$, and $m\in\mathbb N_{>0}$ the worker limit. Ignoring communication and scheduling overhead, the makespan $T$ satisfies
\begin{equation}
 T\geq\max\left\{\frac{\sum_{v\in V}w_v}{m},
 \max_{p\in\mathcal P}\sum_{v\in p}w_v\right\},
\end{equation}
where $\mathcal P$ is the set of dependency paths. The first term is a work bound; the second is the critical-path bound. Adding workers cannot overcome a serial evidence dependency. Communication, verifier queues, tool contention, and failed retries increase observed makespan above this idealized lower bound.

The graph itself may be uncertain. A search task can reveal a missing prerequisite, invalidate a planned branch, or require another tool. Dynamic orchestration therefore operates over both task actions and graph-editing actions. A manager that repeatedly replans can spend more time maintaining the graph than the workers spend executing it. Graph-edit operations should have budgets and stopping rules, particularly when a model is rewarded for producing elaborate plans rather than verified progress.

Decomposition also changes statistical difficulty. Breaking a coherent task into subtasks may remove context that made the original problem solvable. Conversely, it can reduce distraction by exposing each worker to a narrower evidence set. A controlled experiment should distinguish these effects by comparing the same decomposition executed serially, in parallel without communication, and in parallel with communication. If only the communicating condition improves, evidence exchange is implicated. If serial decomposition already captures the gain, parallelism may primarily affect latency.

An important operating condition is interface stability. Delegation works best when a subtask produces a bounded artifact that can be checked independently of the worker's narrative. A parsed dataset, executable patch, or evidence table has a clearer acceptance surface than a general recommendation. This does not make artifact verification trivial, but it constrains where errors can enter integration. Tasks requiring continuous negotiation of a shared latent interpretation are harder to decompose safely than tasks with explicit data dependencies.

The graph also needs an invalidation relation. If node $v$ consumed artifact digest $d$ from node $u$, a later replacement of that artifact must mark $v$ and its affected descendants stale; graph reachability alone is insufficient when several artifact versions coexist. Reexecuting every descendant is conservative but potentially expensive. Reusing a descendant is justified only when its declared dependencies remain valid, or when an independently checked equivalence permits substitution. A model's statement that a revision is minor is not such an equivalence. This turns the critical path into a versioned evidence path: parallel work reduces elapsed computation only to the extent that its outputs survive integration without invalidating subsequent work. Chapter~\ref{ch:C32} applies the same dependency principle to deployed model, tool and memory updates.

